\subsection{A hole does not determine the nearby limit}

Consider the function that agrees with \(y=x+1\) near \(x=1\) but is not defined at \(x=1\). The graph has a hole at \((1,2)\), yet every sequence of inputs approaching \(1\) produces output values approaching \(2\).

\begin{center}
\begin{tikzpicture}[x=1.25cm,y=1.05cm]
  \draw[->] (-2.3,0) -- (3.1,0) node[right] {$x$};
  \draw[->] (0,-1.1) -- (0,4.3) node[above] {$y$};
  \draw[thick,domain=-1.8:2.6,samples=2] plot (\x,{\x+1});
  \fill[white] (1,2) circle (3pt);
  \draw[thick] (1,2) circle (3pt);
  \draw[dashed] (1,0) -- (1,2);
  \draw[dashed] (0,2) -- (1,2);
  \node[below] at (1,0) {$1$};
  \node[left] at (0,2) {$2$};
  \draw[->] (0.15,1.55) -- (0.82,1.93);
  \draw[->] (1.85,2.85) -- (1.18,2.12);
  \node at (-0.3,1.25) {from the left};
  \node at (2.2,3.15) {from the right};
\end{tikzpicture}
\end{center}

\begin{interpretationbox}[Sequential reading of the graph]
Points may approach \(x=1\) from either side or alternate between the two sides. In every case their heights on the line approach \(2\). This is the geometric content of Heine's theorem in this example.
\end{interpretationbox}